%%%%%%%%%%%%%%%%%%%%%%%%%%%%%%%%%%%%%%%%%%%%%%%%%%%%%%%%%%%%%%%%%%%%%%%%%%%
% Lehreinheit 3: Literaturrecherche
% Quelle: WuZ-Buch, Kap. 15 (LE 3), Kap. 5b (Congruence Bias,
%         sec:congruence_bias), Kap. 11 (Effekttabelle: Kruskal Count),
%         Kap. 15b (Kruskal Count als Referenzerfahrung)
% Staunmoment: Kruskal Count  |  Kernbegriff: Suchstrategie
% Der Text fuer den Kruskal Count steht in handouts/LE03.tex und wird mit
% tools/kruskal_check.py geprueft (Schlusswort: "Quelle").
%%%%%%%%%%%%%%%%%%%%%%%%%%%%%%%%%%%%%%%%%%%%%%%%%%%%%%%%%%%%%%%%%%%%%%%%%%%
\nsl{Mit einer Forschungsfrage aus LE~2 beginnt die Suche nach dem, was
andere dazu schon herausgefunden haben. Die meisten Studierenden kennen
dafür ein einziges Werkzeug: das Suchfeld einer Suchmaschine. Diese Einheit
ersetzt es durch eine Strategie. Der Staunmoment, der Kruskal Count, lässt
die ganze Gruppe von verschiedenen Startpunkten aus auf demselben Wort
landen. Das Prinzip dahinter ist mathematisch und darf vollständig erklärt
werden. Es liefert ein Bild für zwei Dinge zugleich: Eine gute Suchstrategie
führt aus vielen Richtungen zu den zentralen Quellen. Eine einseitige
Strategie führt aus vielen Richtungen zur eigenen Vorannahme. Der zweite
Teil dieses Bildes verbindet die Einheit mit dem Bestätigungsfehler, den
das Buch in Kapitel~5b am Wason-Test entfaltet.
\newpage}
\os{\begin{description}
\item[{\rot\bf Heute}]~\\[-\lblskip]
\bi
  \ite ein Text, an dessen Ende alle dasselbe Wort finden
  \ite systematisch statt zufällig suchen
  \ite Datenbanken, boolesche Operatoren, Trunkierung
  \ite Schneeball und Zitationsnetzwerke
  \ite dokumentieren, verwalten, annotieren
  \ite die Pflicht zur Gegenrecherche
\ei
\end{description}
\newpage}

\begin{vorbereitung}[Kruskal Count]
\textbf{Eine Woche vorher}
\begin{checkliste}
  \item Den Text auf dem Übungsblatt Teil~A (Aufgabe~1) selbst dreimal mit
    verschiedenen Startzahlen durchzählen, bis die Regel sicher sitzt. Das
    Schlusswort für alle Startzahlen von 1 bis 10 ist \textbf{\enquote{Quelle}}
    (letztes Wort des Textes). Das Prüfskript
    \texttt{tools/kruskal\_check.py} bestätigt das; wer den Text ändert,
    muss es erneut laufen lassen.
  \item Eine \textbf{Vorhersage} vorbereiten: \enquote{Quelle} auf eine
    Karte schreiben, in einen Umschlag stecken, zukleben und mit Datum
    beschriften. Der Umschlag liegt von Beginn an sichtbar auf dem Pult.
  \item \textbf{Live-Demo planen:} bei der Hochschulbibliothek klären, welche
    Datenbanken lizenziert sind (Web of Science, Scopus, Fachdatenbanken)
    und ob sie im Hörsaal-WLAN ohne VPN erreichbar sind. Eine Beispielfrage
    festlegen (Vorschlag: \enquote{Fördert ein Zaubereffekt im Unterricht
    das Behalten?}) und die Suche einmal komplett durchspielen.
  \item Von der Probesuche Bildschirmfotos der Trefferlisten sichern (Plan B
    bei Netzproblemen). Trefferzahlen \emph{nicht} auf Folien festschreiben;
    sie ändern sich laufend und werden in der Stunde live abgelesen.
  \item Optional: Fachreferent:in der Bibliothek für 10~Minuten in die
    Live-Demo einladen.
\end{checkliste}

\textbf{Handout und Übungsblatt}
\begin{checkliste}
  \item Übungsblatt LE~3 Teil~A und Teil~B je einmal pro Person plus fünf
    Reserve kopieren. Teil~A enthält nur Aufgabe~1 mit dem Text für den
    Staunmoment und wird zu Beginn \emph{verdeckt} ausgeteilt. Den Text
    einseitig und gut lesbar drucken (nicht verkleinern). Teil~B
    (Aufgaben 2--6) wird zusammen mit dem Handout nach der Auflösung
    ausgeteilt.
  \item Handout \enquote{Literaturrecherche} je einmal pro Person; es wird
    erst nach der Auflösung ausgeteilt, weil es das Prinzip erklärt.
  \item Bleistifte bereitlegen; wer mitzählen will, tippt mit dem Stift von
    Wort zu Wort.
\end{checkliste}

\textbf{Am Tag selbst}
\begin{checkliste}
  \item Beamer mit Browser und eingeloggtem Bibliothekszugang testen;
    Literaturverwaltung (Zotero) auf dem Dozierendenrechner geöffnet.
  \item Die Folie mit der Zählregel und dem Übungsbeispiel bereithalten;
    die Regel ist der kritische Punkt.
  \item Tafel: links \enquote{Wie ist das möglich?}, rechts später
    \enquote{Was heißt das für die Suche?}
\end{checkliste}

\textbf{Risiken und Plan B.} Die häufigste Störung sind Zählfehler,
besonders beim ersten Schritt (Wird das Startwort mitgezählt?) und bei
Wörtern mit Umlauten. Deshalb die Regel an einem fremden Beispielsatz
vorführen. Wer auf einem anderen Wort landet, ist kein Ausfall, sondern ein
Datenpunkt: Zwei Personen mit verschiedenen Schlusswörtern zählen
gemeinsam nach und finden den Fehler. Die Kopplung ist robust; ein
einzelner Zählfehler führt oft wieder in den gemeinsamen Weg zurück. Das
Prinzip ist öffentlich und wird vollständig aufgelöst. Fällt das Netz aus,
läuft die Live-Demo mit den gesicherten Bildschirmfotos.
\end{vorbereitung}

\begin{ablauf}
\begin{description}
\item[0--5 · Blitzlicht.] Kurzer Rückblick: Welche Frage haben wir letzte
  Woche geschärft? Danach richtet sich heute die Suche.
\item[5--15 · Staunmoment.] Kruskal Count mit der ganzen Gruppe;
  Vorhersage öffnen; Staunen über die Konvergenz; Erklärungsvermutungen
  sammeln, dann auflösen. Danach Handout und Übungsblatt Teil~B austeilen.
\item[15--45 · Theorie-Input und Live-Demo.] Systematische Recherche,
  Datenbanken, boolesche Operatoren, Schneeball, Literaturverwaltung;
  Live-Recherche an der Beispielfrage. Einschub: Bestätigungsfehler in der
  Suche und die Pflicht zur Gegenrecherche (Buch, Kap.~5b).
\item[45--75 · Rechercheprotokoll.] Die Studierenden starten die
  dokumentierte Recherche zur eigenen Frage (Übungsblatt Teil~B,
  Aufgaben 2--4).
\item[75--90 · Sicherung und Transfer.] Erste Treffer nach Stufen sortieren
  (Aufgabe~5); Auftrag zur annotierten Bibliographie als Brücke zu LE~4
  (Quellenkritik).
\end{description}
\end{ablauf}

%==============================================================================
\section{Staunmoment: Der Zählweg}
%==============================================================================
\begin{staunbox}{Kruskal Count}
Wählen Sie im Stillen eine Zahl von 1 bis 10. Zählen Sie im Text auf dem
Übungsblatt so viele Wörter ab: Das ist Ihr Startwort.

\medskip
Zählen Sie seine Buchstaben und gehen Sie so viele Wörter weiter. Wieder
Buchstaben zählen, wieder weitergehen. Wenn der nächste Sprung über das
Textende hinausführt, haben Sie Ihr \textbf{Schlusswort}.

\medskip
\emph{Behalten Sie es für sich.}
\end{staunbox}
\os{\newpage}

\begin{description}
\item[{\rot\bf Die Regel an einem Beispiel}]~\\[-\lblskip]
\bi
  \ite Satz: \enquote{Der kleine Hund sprang über den breiten Graben und
    bellte laut.}
  \ite Startzahl 2 $\to$ Startwort \enquote{kleine} (6 Buchstaben)
  \ite 6 Wörter weiter: Hund, sprang, über, den, breiten, \textbf{Graben}
  \ite \enquote{Graben} hat 6 Buchstaben, es folgen nur noch 3 Wörter
    $\to$ Schlusswort \enquote{Graben}
  \ite Satzzeichen zählen nicht; ä, ö, ü zählen als ein Buchstabe
\ei
\nsl{Das Beispiel stammt nicht aus dem Übungstext, damit niemand vorab
etwas über das Schlusswort erfährt. Wichtig ist der erste Sprung: Man
zählt ab dem \emph{nächsten} Wort. Nach dem Beispiel zählen alle still.
Wer fertig ist, dreht das Blatt um. Nach etwa acht Minuten fragt die
Lehrperson, wer ein Schlusswort hat, und lässt eine Person den Umschlag
öffnen.}
\os{\newpage}

\item[{\rot\bf Was ist gerade passiert?}]~\\[-\lblskip]
\bi
  \ite Wer hat \enquote{Quelle} als Schlusswort? Wer nicht?
  \ite Sie haben Ihre Zahl frei gewählt. Wie konnte die Vorhersage
    stimmen?
  \ite sammeln Sie Erklärungen: Zufall? ein präparierter Text? etwas
    anderes?
  \ite Wie könnte man die Erklärungen prüfen?
\ei
\nsl{Typische Vermutungen: Der Text sei so gebaut, dass alle Wege
konstruiert zusammenlaufen; es gebe nur ein einziges Wort, das man
erreichen kann; die Lehrperson habe die Wahl beeinflusst. Die ersten beiden
liegen nahe an der Wahrheit, die dritte ist hier falsch, aber eine gute
Frage: Im Buch dient der Kruskal Count als Beispiel für eine scheinbar
freie Wahl, deren Ergebnis feststeht (Kap.~15b). Prüfbar wird das, wenn
zwei Personen mit verschiedenen Startzahlen ihre Wege nebeneinanderlegen
und nachsehen, wo sie sich treffen.}
\os{\newpage}
\end{description}

%==============================================================================
\section{Auflösung: das Kruskal-Prinzip}
%==============================================================================
\nsl{Der Effekt geht auf den Mathematiker und Physiker Martin Kruskal zurück
und wurde durch Martin Gardner bekannt; eine mathematische Analyse geben
\textcite{lagarias2009kruskal}. Die Erklärung beruht auf einer einfachen
Beobachtung: Der Weg durch den Text hängt nur vom aktuellen Wort ab, nicht
davon, wie man dorthin gekommen ist. Treffen zwei Wege einmal auf
dasselbe Wort, gehen sie von da an denselben Weg bis zum Ende. In der
Sprache der Wahrscheinlichkeitstheorie sind das zwei Markov-Ketten, die
\emph{koppeln}.

Wie wahrscheinlich ist ein solches Treffen? Deutsche Wörter haben im
Mittel einige Buchstaben; jeder Weg trifft also ungefähr jedes vierte bis
sechste Wort. Zwei Wege, die nebeneinander herlaufen, haben bei jedem
Schritt eine gute Chance, auf dasselbe Wort zu fallen. Über viele Schritte
wird es fast sicher, dass sie es irgendwann tun. Im Übungstext laufen alle
zehn Startwege nach dem Prüfskript spätestens beim 47.~Wort zusammen; die
übrigen rund 260 Wörter legen sie gemeinsam zurück. Ein langer Text ist
also kein Hindernis für den Effekt, er ist seine Voraussetzung.}
\begin{description}
\item[{\rot\bf Zwei Wege, die sich treffen}]~\\[-\lblskip]
\begin{center}
\resizebox{0.95\textwidth}{!}{%
\begin{tikzpicture}[font=\sffamily\small,>=Stealth]
  \foreach \i in {1,...,22} {\draw[rwugray!60] (\i*0.6,0) rectangle +(0.5,0.5);}
  % Weg A (violett) ueber der Leiste
  \foreach \a/\b in {1/4,4/9,9/12,12/17,17/21} {
    \draw[->,thick,rwuviolet] (\a*0.6+0.25,0.55) to[bend left=50] (\b*0.6+0.25,0.55);}
  % Weg B (cyan) unter der Leiste
  \foreach \a/\b in {2/7,7/12} {
    \draw[->,thick,rwucyan] (\a*0.6+0.25,-0.05) to[bend right=50] (\b*0.6+0.25,-0.05);}
  \fill[rwuviolet!40] (12*0.6,0) rectangle +(0.5,0.5);
  \node[below=6mm,align=center] at (12*0.6+0.25,0) {Treffpunkt:\\ab hier ein Weg};
  \node[rwuviolet,left] at (0.5,0.9) {Start 1};
  \node[rwucyan,left] at (0.5,-0.4) {Start 2};
\end{tikzpicture}}
\end{center}
\bi
  \ite der nächste Schritt hängt nur vom aktuellen Wort ab
  \ite einmal getroffen, immer gemeinsam: Die Wege \emph{koppeln}
\ei
\os{\newpage}

\item[{\rot\bf Warum fast alle treffen}]~\\[-\lblskip]
\bi
  \ite jeder Weg berührt etwa jedes vierte bis sechste Wort
  \ite bei jedem Schritt besteht eine Chance, auf ein Wort eines anderen
    Wegs zu fallen
  \ite über viele Schritte wird das Treffen fast sicher
  \ite im Übungstext: alle zehn Wege treffen sich spätestens bei Wort 47
    von 312
  \ite Analyse: \textcite{lagarias2009kruskal}
\ei
\os{\newpage}

\item[{\rot\bf Was der Trick über Suchen sagt}]~\\[-\lblskip]
\bi
  \ite \textbf{Konvergenz:} Eine gute Strategie führt von vielen
    Startpunkten zu den zentralen Quellen
  \ite \textbf{Kehrseite:} Eine einseitige Strategie führt von vielen
    Startpunkten zur eigenen Vorannahme
  \ite das Ergebnis kann feststehen, obwohl sich jeder Schritt frei
    anfühlt
  \ite Prüffrage aus dem Buch: \enquote{Ist das gerade ein Kruskal Count?}
\ei
\nsl{Die Prüffrage stammt aus Kapitel~15b des Buchs, das den Trick als
gemeinsame Referenzerfahrung für Teams vorschlägt: Wer ihn erlebt hat, kann
fragen, ob ein Ergebnis wirklich aus der Sache folgt oder aus der Regel, nach
der man sucht. Dass ein Weg sich frei anfühlt, sagt wenig darüber, wohin er
führt.}
\os{\newpage}
\end{description}

%==============================================================================
\section{Systematisch statt zufällig}
%==============================================================================
\nsl{Eine unsystematische Recherche beginnt mit einem Suchbegriff, folgt
dem ersten interessanten Treffer und endet, wenn die Zeit knapp wird oder
man genug gefunden zu haben glaubt. Das Ergebnis hängt dann vom
Startpunkt, von der Reihenfolge der Trefferliste und von der Laune des
Tages ab. Niemand, auch man selbst nicht, kann es später nachvollziehen. Eine
systematische Recherche legt dagegen vorher fest, mit welchen Begriffen,
in welchen Quellen und nach welchen Kriterien gesucht wird, und
protokolliert jeden Schritt \parencite{doering2016forschungsmethoden}.
In ihrer strengsten Form, der systematischen Übersichtsarbeit, wird sogar
das Protokoll vorab veröffentlicht und die Auswahl in einem Flussdiagramm
offengelegt \parencite{page2021prisma}. Für eine Seminar- oder
Bachelorarbeit genügt eine einfachere Fassung. Der Grundsatz bleibt
derselbe: Die Suche muss sich wiederholen lassen.}
\begin{description}
\item[{\rot\bf Zwei Arten zu suchen}]~\\[-\lblskip]
\begin{center}
\small
\begin{tabular}{p{0.42\textwidth}p{0.42\textwidth}}
\toprule
\textbf{unsystematisch} & \textbf{systematisch} \\
\midrule
ein Suchbegriff, eine Suchmaschine & Begriffe pro Konzept, mehrere Datenbanken \\
dem ersten guten Treffer folgen & Ein- und Ausschlusskriterien vorab \\
aufhören, wenn es reicht & aufhören, wenn nichts Neues mehr kommt \\
nicht nachvollziehbar & Protokoll: wiederholbar \\
\bottomrule
\end{tabular}
\end{center}
\os{\newpage}

\item[{\rot\bf Arbeitsdefinition}]~\\
\defn{systrecherche}{Eine \textbf{systematische Literaturrecherche} ist
eine geplante Suche, bei der Suchbegriffe, Quellen und Auswahlkriterien
vorab festgelegt und alle Schritte so dokumentiert werden, dass eine andere
Person die Suche wiederholen und zum selben Ergebnis kommen kann.}
\bi
  \ite \emph{geplant}: Begriffe und Kriterien stehen vor der Suche fest
  \ite \emph{dokumentiert}: Rechercheprotokoll
  \ite \emph{reproduzierbar}: der Prüfstein jeder Recherche
\ei
\os{\newpage}

\item[{\rot\bf Der Ablauf in fünf Schritten}]~\\[-\lblskip]
\begin{center}
\resizebox{0.95\textwidth}{!}{%
\begin{tikzpicture}[font=\sffamily\small,>=Stealth,
  st/.style={draw=rwuviolet,thick,fill=rwuviolet!8,rounded corners=2pt,
    minimum height=18mm,align=center,text width=26mm}]
  \node[st] (a) at (0,0) {Frage in Konzepte zerlegen};
  \node[st] (b) at (3.2,0) {Begriffe und Synonyme sammeln};
  \node[st] (c) at (6.4,0) {Suchstring in Datenbanken};
  \node[st] (d) at (9.6,0) {Schneeball rückwärts und vorwärts};
  \node[st,draw=rwucyan,fill=rwucyan40!40] (e) at (12.8,0) {sichten, auswählen, protokollieren};
  \draw[->,thick] (a) -- (b); \draw[->,thick] (b) -- (c);
  \draw[->,thick] (c) -- (d); \draw[->,thick] (d) -- (e);
  \draw[->,thick,rwugray,dashed] (e.south) to[bend left=15]
    node[below]{neue Begriffe aus den Treffern $\to$ nächste Runde} (b.south);
\end{tikzpicture}}
\end{center}
\os{\newpage}
\end{description}

%==============================================================================
\section{Wo man sucht: Kataloge und Datenbanken}
%==============================================================================
\nsl{Ein Bibliothekskatalog verzeichnet, was eine Bibliothek besitzt, vor
allem Bücher und Zeitschriftentitel, aber nicht die einzelnen Aufsätze darin.
Für Aufsätze braucht man Datenbanken. Fachübergreifende
Zitationsdatenbanken wie Web of Science und Scopus erfassen ausgewählte
Zeitschriften mit ihren Literaturlisten und erlauben deshalb auch die Suche
nach zitierenden Arbeiten. Fachdatenbanken (PubMed für Medizin, PsycInfo
für Psychologie, ERIC und das Fachportal Pädagogik für Bildung, IEEE Xplore
und die ACM Digital Library für Technik und Informatik) erschließen ein Fach
tiefer, oft mit einem kontrollierten Schlagwortverzeichnis. Google Scholar
findet sehr viel, ist aber für eine systematische Suche nur bedingt
geeignet: Die Suchsyntax ist eingeschränkt, Trefferlisten lassen sich nicht
exakt wiederholen, und es bleibt unklar, was überhaupt erfasst ist.
\textcite{gusenbauer2020academic} haben 28 Suchsysteme verglichen und
empfehlen Google Scholar deshalb nur als Ergänzung, nicht als Hauptquelle
einer systematischen Suche.}
\begin{description}
\item[{\rot\bf Werkzeuge und wofür sie taugen}]~\\[-\lblskip]
\begin{center}
\small
\begin{tabular}{p{0.32\textwidth}p{0.55\textwidth}}
\toprule
\textbf{Werkzeug} & \textbf{wofür} \\
\midrule
Bibliothekskatalog, KVK & Bücher, Zeitschriftentitel, Bestand vor Ort \\
Web of Science, Scopus & fachübergreifend, mit Zitationssuche \\
PubMed, PsycInfo, ERIC, IEEE Xplore \dots & Fachdatenbanken mit Schlagwörtern \\
Google Scholar & breiter Einstieg, Ergänzung, Volltexte finden \\
BASE, Repositorien & Open Access, graue Literatur \\
\bottomrule
\end{tabular}
\end{center}
\os{\newpage}

\item[{\rot\bf Google Scholar: ja, aber nicht allein}]~\\[-\lblskip]
\bi
  \ite findet viel, auch Volltexte und graue Literatur
  \ite eingeschränkte Suchsyntax, keine Trunkierung
  \ite Trefferlisten nicht exakt wiederholbar
  \ite unklar, was erfasst ist und was nicht
  \ite Empfehlung: Ergänzung, nicht Hauptquelle
    \parencite{gusenbauer2020academic}
\ei
\os{\newpage}
\end{description}

%==============================================================================
\section{Suchbegriffe und Suchsprache}
%==============================================================================
\nsl{Datenbanken finden nur, was man ihnen sagt. Der erste Schritt ist
deshalb, die Forschungsfrage in zwei bis drei Konzepte zu zerlegen und zu
jedem Konzept die Begriffe zu sammeln, unter denen Fachleute darüber
schreiben: Synonyme, englische Fachbegriffe, Ober- und Unterbegriffe. Viele
Fachdatenbanken bieten dafür einen \emph{Thesaurus}, ein kontrolliertes
Schlagwortverzeichnis (in PubMed etwa die Medical Subject Headings). Wer
nach dem Schlagwort sucht statt nach dem Wort im Titel, findet auch
Arbeiten, die andere Formulierungen verwenden.

Die Suchbegriffe verbindet man mit booleschen Operatoren. Innerhalb eines
Konzepts stehen die Synonyme mit \texttt{OR}, denn jedes genügt;
zwischen den Konzepten steht \texttt{AND}, denn alle müssen vorkommen.
\texttt{NOT} schließt aus und ist gefährlich, weil es auch relevante Arbeiten
entfernt, in denen der ausgeschlossene Begriff nur am Rand vorkommt.
Trunkierung (meist mit \texttt{*}) findet verschiedene Wortendungen,
Anführungszeichen erzwingen eine feste Wortfolge.}
\begin{description}
\item[{\rot\bf Boolesche Operatoren}]~\\[-\lblskip]
\begin{center}
\resizebox{0.85\textwidth}{!}{%
\begin{tikzpicture}[font=\sffamily]
  \begin{scope}
    \clip (0,0) circle (0.9);
    \fill[rwuviolet!45] (1.1,0) circle (0.9);
  \end{scope}
  \draw[thick] (0,0) circle (0.9) (1.1,0) circle (0.9);
  \node at (0.55,-1.3) {\texttt{A AND B}};
  \node at (0.55,-1.8) {\footnotesize weniger, genauer};
  \begin{scope}[xshift=4cm]
    \fill[rwucyan40!70] (0,0) circle (0.9) (1.1,0) circle (0.9);
    \draw[thick] (0,0) circle (0.9) (1.1,0) circle (0.9);
    \node at (0.55,-1.3) {\texttt{A OR B}};
    \node at (0.55,-1.8) {\footnotesize mehr, breiter};
  \end{scope}
  \begin{scope}[xshift=8cm]
    \fill[rwuviolet!25] (0,0) circle (0.9);
    \fill[white] (1.1,0) circle (0.9);
    \draw[thick] (0,0) circle (0.9) (1.1,0) circle (0.9);
    \node at (0.55,-1.3) {\texttt{A NOT B}};
    \node at (0.55,-1.8) {\footnotesize Vorsicht: verliert Treffer};
  \end{scope}
  \node at (0,0) {A}; \node at (1.1,0) {B};
  \node at (4,0) {A}; \node at (5.1,0) {B};
  \node at (7.8,0) {A}; \node at (9.1,0) {B};
\end{tikzpicture}}
\end{center}
\os{\newpage}

\item[{\rot\bf Werkzeuge der Suchsprache}]~\\[-\lblskip]
\bi
  \ite \textbf{Trunkierung:} \texttt{learn*} findet learn, learning,
    learner
  \ite \textbf{Phrase:} \texttt{\symbol{34}magic trick\symbol{34}} nur als
    feste Wortfolge
  \ite \textbf{Klammern:} \texttt{(magic OR conjuring) AND learn*}
  \ite \textbf{Feldsuche:} nur in Titel, Abstract oder Schlagwort
  \ite \textbf{Thesaurus:} kontrollierte Schlagwörter der Datenbank
  \ite Syntax unterscheidet sich je Datenbank: Hilfeseite lesen
\ei
\os{\newpage}

\item[{\rot\bf Live-Demo: eine Beispielfrage}]~\\[-\lblskip]
\bi
  \ite Frage: Fördert ein Zaubereffekt im Unterricht das Behalten?
  \ite Konzepte: \emph{Zauberkunst} -- \emph{Staunen/Neugier} --
    \emph{Lernen}
  \ite Suchstring:\\
    {\small\ttfamily (magic OR conjuring OR \symbol{34}magic trick*\symbol{34})\\
    AND (wonder OR awe OR curiosity OR surprise)\\
    AND (learn* OR memory OR retention)}
  \ite Was passiert mit der Trefferzahl, wenn wir ein Konzept weglassen?
\ei
\nsl{In der Live-Demo wird derselbe String in zwei Datenbanken und in
Google Scholar eingegeben. Die Studierenden sollen sehen, dass die
Trefferzahlen stark schwanken, dass AND die Menge verkleinert und OR sie
vergrößert, und dass ein einziger zusätzlicher Begriff die Liste deutlich
verändern kann. Die Zahlen werden live abgelesen und nicht vorab auf Folien
geschrieben. Als Startpunkt für die Schneeballsuche eignen sich die
Arbeiten, die das Buch selbst zitiert, etwa zur Neugier
\parencite{loewenstein1994} oder zur Wahrnehmungspsychologie der Zauberkunst.
Ein guter Moment für die Gegenfrage: Mit welchen Begriffen fänden wir
Arbeiten, die \emph{keinen} Lerneffekt zeigen?}
\os{\newpage}
\end{description}

%==============================================================================
\section{Schneeball und Zitationsnetzwerke}
%==============================================================================
\nsl{Die Schlagwortsuche hat eine natürliche Grenze: Sie findet nur
Arbeiten, deren Wortwahl man vorausgeahnt hat. Ein zweiter Weg führt über
die Literaturlisten. Bei der \emph{Rückwärtssuche} liest man die
Literaturliste eines Schlüsseltexts und prüft die dort zitierten Arbeiten.
Bei der \emph{Vorwärtssuche} fragt man, welche späteren Arbeiten den
Schlüsseltext zitieren. Die Vorwärtssuche wurde erst durch Zitationsindizes
möglich, die Eugene Garfield in den 1950er-Jahren vorschlug
\parencite{garfield1955citation} und auf die Web of Science zurückgeht.
\textcite{wohlin2014snowballing} hat daraus ein geregeltes Verfahren
gemacht: Man beginnt mit einer kleinen Startmenge, sucht in Runden
rückwärts und vorwärts und hört auf, wenn keine neuen relevanten Arbeiten
mehr auftauchen.

Hier wird der Kruskal Count zum genauen Bild. Wer von verschiedenen
Schlüsseltexten aus den Literaturlisten folgt, landet nach wenigen Schritten
bei denselben häufig zitierten Arbeiten. Das ist nützlich, denn so findet
man die Klassiker eines Gebiets. Es hat aber denselben Haken wie der Trick:
Die Konvergenz sagt etwas über das Netz, nicht über die Wahrheit. Häufig
zitiert wird auch, was in einer Gemeinschaft gern gelesen wird.}
\begin{description}
\item[{\rot\bf Rückwärts und vorwärts}]~\\[-\lblskip]
\begin{center}
\resizebox{0.9\textwidth}{!}{%
\begin{tikzpicture}[font=\sffamily,>=Stealth,
  p/.style={draw=rwugray,fill=white,rounded corners=1pt,minimum width=14mm,minimum height=8mm},
  k/.style={draw=rwuviolet,very thick,fill=rwuviolet!15,rounded corners=1pt,minimum width=20mm,minimum height=10mm}]
  \node[k] (k) at (0,0) {Schlüsseltext};
  \node[p] (a1) at (-4.5,1.2) {1998};
  \node[p] (a2) at (-4.5,0) {2004};
  \node[p] (a3) at (-4.5,-1.2) {2009};
  \node[p] (b1) at (4.5,1.2) {2019};
  \node[p] (b2) at (4.5,0) {2021};
  \node[p] (b3) at (4.5,-1.2) {2024};
  \foreach \x in {a1,a2,a3} \draw[->,thick,rwuviolet] (k) -- (\x);
  \foreach \x in {b1,b2,b3} \draw[->,thick,rwucyan] (\x) -- (k);
  \node[rwuviolet,align=center] at (-4.5,-2.3) {\textbf{rückwärts}\\Literaturliste};
  \node[rwucyan,align=center] at (4.5,-2.3) {\textbf{vorwärts}\\wer zitiert ihn?};
  \draw[->,thick,rwugray] (-6,-3.2) -- node[below]{Zeit} (6,-3.2);
\end{tikzpicture}}
\end{center}
\os{\newpage}

\item[{\rot\bf Schneeball als Verfahren}]~\\[-\lblskip]
\bi
  \ite Startmenge: zwei bis fünf sicher relevante Texte
  \ite Runde für Runde: rückwärts und vorwärts
  \ite Einschlusskriterien wie bei der Schlagwortsuche
  \ite Stopp, wenn keine neuen relevanten Treffer kommen
    \parencite{wohlin2014snowballing}
  \ite Vorwärtssuche: Web of Science, Scopus, Google Scholar
    (\enquote{Zitiert von})
\ei
\os{\newpage}

\item[{\rot\bf Konvergenz im Netz}]~\\[-\lblskip]
\bi
  \ite verschiedene Startpunkte führen zu denselben Klassikern
  \ite gut: so findet man den Kern eines Gebiets
  \ite aber: häufig zitiert heißt nicht richtig
  \ite Zitiergemeinschaften können eine Lücke um sich herum haben
  \ite darum Schneeball \emph{und} Schlagwortsuche kombinieren
\ei
\os{\newpage}
\end{description}

%==============================================================================
\section{Graue Literatur}
%==============================================================================
\nsl{Als \emph{graue Literatur} bezeichnet man Texte, die nicht über
Verlage und begutachtete Zeitschriften erscheinen: Abschlussarbeiten,
Preprints, Berichte von Behörden, Stiftungen und Verbänden,
Konferenzfolien, Daten- und Projektberichte. Sie ist für eine Recherche aus
zwei Gründen wichtig. Sie ist oft aktueller als begutachtete Literatur, und
sie enthält häufiger Ergebnisse, die keine Zeitschrift drucken wollte,
etwa Studien ohne signifikanten Effekt. Wer sie ganz weglässt, sieht ein zu
günstiges Bild des Forschungsstands ($\to$~Publikationsbias, LE~10). Ihr
Nachteil ist die fehlende Begutachtung: Ihre Qualität muss man selbst
prüfen. Damit beschäftigt sich LE~4.}
\begin{description}
\item[{\rot\bf Graue Literatur: Chancen und Grenzen}]~\\[-\lblskip]
\bi
  \ite Abschlussarbeiten, Preprints, Berichte, Konferenzbeiträge
  \ite \textbf{Chance:} aktuell, enthält auch Nullbefunde
  \ite \textbf{Grenze:} keine Begutachtung, oft schwer zu finden
  \ite Fundorte: BASE, Repositorien, Preprint-Server, Websites von
    Institutionen
  \ite Qualität selbst prüfen ($\to$~LE~4)
\ei
\os{\newpage}
\end{description}

%==============================================================================
\section{Dokumentieren, verwalten, annotieren}
%==============================================================================
\nsl{Das Rechercheprotokoll ist das Laborbuch der Literaturarbeit. Es hält
für jeden Suchvorgang Datum, Datenbank, genauen Suchstring, gesetzte Filter,
Trefferzahl und die Zahl der gesichteten und ausgewählten Arbeiten fest,
dazu die Ein- und Ausschlusskriterien. Die PRISMA-Leitlinie verlangt für
systematische Übersichtsarbeiten zusätzlich ein Flussdiagramm, das zeigt, wie
aus vielen Treffern wenige eingeschlossene Arbeiten wurden
\parencite{page2021prisma}. Für Seminararbeiten ist das ein gutes Vorbild,
auch wenn man es nicht vollständig übernimmt.

Für die gefundenen Quellen braucht man ein Literaturverwaltungsprogramm.
Zotero und JabRef sind frei verfügbar, Citavi ist an vielen Hochschulen
lizenziert. Sie speichern die bibliographischen Angaben, PDFs und Notizen
und erzeugen Literaturverzeichnisse in jedem Zitierstil. Die annotierte
Bibliographie geht einen Schritt weiter: Zu jeder Quelle steht ein kurzer
Kommentar zu Kernaussage, Methode und Relevanz. Das Begleitbuch dieser
Vorlesung ist selbst so gebaut. Sein Literaturverzeichnis druckt zu jeder
Quelle ein Abstract, und das Skript übernimmt diese Form.}
\begin{description}
\item[{\rot\bf Das Rechercheprotokoll}]~\\[-\lblskip]
\begin{center}
\small
\begin{tabular}{llllll}
\toprule
\textbf{Datum} & \textbf{Datenbank} & \textbf{Suchstring} & \textbf{Filter} & \textbf{Treffer} & \textbf{gewählt} \\
\midrule
\dots & Scopus & \texttt{(magic OR \dots) AND \dots} & ab 2000 & \dots & \dots \\
\dots & PsycInfo & \dots & Englisch & \dots & \dots \\
\bottomrule
\end{tabular}
\end{center}
\bi
  \ite Prüfstein: Könnte jemand anderes die Suche wiederholen?
  \ite Ein- und Ausschlusskriterien vorab festlegen
  \ite Vorbild für große Arbeiten: PRISMA-Flussdiagramm
    \parencite{page2021prisma}
\ei
\os{\newpage}

\item[{\rot\bf Literaturverwaltung}]~\\[-\lblskip]
\bi
  \ite \textbf{Zotero}, \textbf{JabRef}: frei; \textbf{Citavi}: oft
    Hochschullizenz
  \ite Angaben per Browser-Plugin oder DOI übernehmen, nie abtippen
  \ite PDFs, Notizen und Schlagwörter an der Quelle ablegen
  \ite Export nach BibTeX/biblatex oder Word
  \ite ab der ersten Quelle verwenden, nicht erst beim Schreiben
\ei
\os{\newpage}

\item[{\rot\bf Annotierte Bibliographie}]~\\[-\lblskip]
\bi
  \ite je Quelle zwei bis drei Sätze
  \ite \textbf{Kernaussage:} Was wird behauptet oder gefunden?
  \ite \textbf{Methode:} Wie wurde das gezeigt?
  \ite \textbf{Relevanz:} Was bedeutet das für meine Frage?
  \ite Vorbild: das Literaturverzeichnis des Begleitbuchs, mit Abstract zu
    jeder Quelle
\ei
\os{\newpage}
\end{description}

%==============================================================================
\section{Die Kehrseite: suchen, was man schon glaubt}
%==============================================================================
\nsl{Der Kruskal Count hat gezeigt, dass eine feste Regel aus vielen
Startpunkten zum selben Ziel führt. Bei der Recherche ist die Regel oft die
eigene Erwartung. Wer überzeugt ist, dass Zauberkunst das Lernen fördert,
sucht nach \texttt{magic AND learning AND benefit}, liest die passenden
Treffer zuerst und folgt deren Literaturlisten. So entsteht ein Netz von
Bestätigungen, das sich wie Forschungsstand anfühlt.

Das Buch beschreibt in Kapitel~5b den psychologischen Kern dieses Musters
am 2-4-6-Test von \textcite{wason1960failure}. Die Versuchspersonen sollten
eine Regel hinter dem Tripel 2--4--6 erraten und durften dazu eigene Tripel
testen. Die Regel war weit gefasst (drei aufsteigende Zahlen), doch die
meisten legten sich auf eine engere Vermutung fest und prüften nur noch
Tripel, die zu ihr passten. Nur sechs von 29 Personen fanden die Regel ohne
falschen Zwischenschluss. Jonathan Baron nennt die Neigung, nur die eigene
Hypothese direkt zu prüfen, den \emph{Congruence Bias} und formuliert als
Gegenmittel: an Alternativhypothesen denken und den Test wählen, der
zwischen ihnen unterscheidet \parencite{baron1988diagnostic,baron2008thinking}.
\textcite{klaymanha1987} haben gezeigt, dass diese positive Teststrategie in
vielen Situationen vernünftig ist und erst dann in die Irre führt, wenn die
Wahrheit breiter ist als die eigene Vermutung. Dass derselbe Mechanismus in
Wissenschaft, Medizin und Justiz allgegenwärtig ist, dokumentiert
\textcite{nickerson1998ubiquitous}. Für die Recherche folgt daraus eine
Pflicht: Man sucht gezielt auch nach Arbeiten, die der eigenen These
widersprechen könnten.}
\begin{description}
\item[{\rot\bf Der 2-4-6-Test}]~\\[-\lblskip]
\bi
  \ite Aufgabe: die Regel hinter 2--4--6 finden, eigene Tripel testen
  \ite Regel: drei aufsteigende Zahlen
  \ite die meisten testen nur Tripel, die zu \emph{ihrer} Vermutung passen
    (8--10--12, 20--22--24 \dots)
  \ite nur sechs von 29 fanden die Regel ohne falschen Zwischenschluss
    \parencite{wason1960failure}
\ei
\os{\newpage}

\item[{\rot\bf Congruence Bias in der Recherche}]~\\[-\lblskip]
\bi
  \ite Suchbegriffe enthalten schon die erwartete Antwort
  \ite passende Treffer werden zuerst gelesen
  \ite Schneeball folgt den Literaturlisten der Bestätigungen
  \ite Ergebnis: ein Netz, das sich wie Forschungsstand anfühlt
  \ite Gegenmittel: Alternativen denken, unterscheidende Suche wählen
    \parencite{baron2008thinking}
\ei
\os{\newpage}

\item[{\rot\bf Gegenrecherche: so geht's}]~\\[-\lblskip]
\bi
  \ite die Gegenhypothese als eigenes Konzept suchen
  \ite Begriffe für Nullbefunde und Kritik:\\
    \texttt{\symbol{34}no effect\symbol{34}},
    \texttt{\symbol{34}null result*\symbol{34}},\\
    \texttt{replicat*}, \texttt{critique}
  \ite in der Vorwärtssuche gezielt kritische Repliken lesen
  \ite im Protokoll festhalten, was die Gegenrecherche ergab
\ei
\os{\newpage}
\end{description}

\begin{merkbox}
Eine Recherche ist erst dann wissenschaftlich, wenn jemand anderes sie mit
Ihrem Protokoll wiederholen könnte und wenn Sie auch nach dem gesucht
haben, was Ihnen widerspricht. Dass viele Wege zur selben Quelle führen,
spricht für eine gute Suche, beweist aber noch nicht, dass die Quelle recht
hat.
\end{merkbox}
\os{\newpage}

%==============================================================================
\section{Übung in der Stunde}
%==============================================================================
\begin{description}
\item[{\rot\bf Rechercheprotokoll starten (30 Min.)}]~\\[-\lblskip]
\bi
  \ite Zerlegen Sie Ihre Forschungsfrage in Konzepte, sammeln Sie Begriffe
    (Aufgabe~2)
  \ite Bauen Sie einen Suchstring und einen Gegen-Suchstring (Aufgabe~3)
  \ite Suchen Sie in zwei Datenbanken, protokollieren Sie jede Suche
    (Aufgabe~4)
  \ite Wählen Sie einen Schlüsseltext für die Schneeballsuche
\ei
\os{\newpage}

\item[{\rot\bf Sicherung: Treffer sortieren (15 Min.)}]~\\[-\lblskip]
\bi
  \ite Ordnen Sie Ihre ersten fünf Treffer den Stufen A--E zu
    (Aufgabe~5)
  \ite Welche Stufe dominiert? Woran liegt das?
  \ite Hat die Gegenrecherche etwas Neues gebracht?
\ei
\os{\newpage}
\end{description}

\begin{lehrhinweis}
In der Protokollphase sind die häufigsten Fehler: nur Google Scholar, nur
deutsche Begriffe, \texttt{AND} zwischen Synonymen, Suchstrings ohne
Klammern und ein Protokoll, in dem der genaue String fehlt. Die Lehrperson
geht herum und bittet einzelne Studierende, ihre Suche nach dem Protokoll
eines Nachbarn zu wiederholen: Kommt dieselbe Trefferzahl heraus? Für die
Sortierung gilt: Die Stufen A--E sind eine grobe Vorsortierung nach
Publikationstyp. Ein begutachteter Artikel kann schwach sein und ein
Behördenbericht stark. Diese Unterscheidung ist Gegenstand von LE~4 und
sollte hier nur angekündigt werden. Hat jemand nach 30~Minuten keine
passenden Treffer, ist die Frage oft zu eng oder zu ungewöhnlich formuliert;
dann ein Konzept weglassen oder einen Oberbegriff wählen.
\end{lehrhinweis}

%==============================================================================
\section{Lernziele und Aufgaben}
%==============================================================================
\begin{description}
\item[{\rot\bf Nach dieser Einheit können Sie}]~\\[-\lblskip]
\bi
  \ite das Kruskal-Prinzip erklären und auf Suchstrategien übertragen
  \ite eine systematische Recherche planen und dokumentieren
  \ite Suchstrings mit booleschen Operatoren und Trunkierung bauen
  \ite die wichtigsten Datenbanken Ihres Fachs nennen
  \ite Treffer nach Relevanz und Publikationstyp vorsortieren
  \ite eine Gegenrecherche durchführen
\ei
\os{\newpage}

\item[{\rot\bf Bis zur nächsten Woche}]~\\[-\lblskip]
\bi
  \ite \textbf{Rechercheprotokoll} vervollständigen: Suchbegriffe,
    Datenbanken, Treffermengen, Auswahlkriterien (10\,\% der
    Portfolionote)
  \ite \textbf{Annotierte Bibliographie:} fünf Quellen mit je zwei bis drei
    Sätzen zu Kernaussage, Methode und Relevanz; mindestens eine aus der
    Gegenrecherche
  \ite \textbf{Tagebuch des Staunens:} Welcher Fund hat Sie bei der
    Recherche überrascht?
  \ite die fünf Quellen zu LE~4 mitbringen: Wir prüfen ihre Qualität
\ei
\os{\newpage}
\end{description}

\handoutinput{handouts/LE03}
